\documentclass[11pt,a4paper]{article}
\usepackage[T1]{fontenc}
\usepackage[utf8]{inputenc}
\usepackage[czech]{babel}
\usepackage{lmodern,microtype}
\usepackage{amsmath,amssymb,mathtools}
\usepackage{geometry}
\geometry{left=19mm,right=19mm,top=21mm,bottom=21mm}
\usepackage[dvipsnames]{xcolor}
\usepackage{enumitem,tabularx,booktabs}
\usepackage[most]{tcolorbox}
\usepackage{fancyhdr,lastpage}
\usepackage[hidelinks]{hyperref}
\definecolor{deep}{HTML}{203955}
\definecolor{teal}{HTML}{176E76}
\definecolor{slate}{HTML}{506173}
\definecolor{pale}{HTML}{F3F7F9}
\setlength{\parindent}{0pt}
\setlength{\parskip}{5pt}
\setlist[enumerate]{leftmargin=6mm,itemsep=5pt,topsep=3pt}
\pagestyle{fancy}\fancyhf{}
\fancyhead[L]{\sffamily\footnotesize\color{slate} MAT · 4. ročník · kombinatorika}
\fancyhead[R]{\sffamily\footnotesize\color{slate} Hodina \LESSONNUM\, / 112}
\fancyfoot[L]{\sffamily\footnotesize\color{slate} Výukový podklad · samostatný dokument}
\fancyfoot[R]{\sffamily\footnotesize\color{slate} \thepage/\pageref{LastPage}}
\renewcommand{\headrulewidth}{0.35pt}
\renewcommand{\footrulewidth}{0pt}
\newcommand{\sect}[1]{\par\vspace{2mm}{\Large\sffamily\bfseries\color{deep} #1}\par\vspace{-0.5mm}\noindent\color{teal}\rule{\linewidth}{0.7pt}\color{black}\vspace{1mm}}
\tcbset{theobox/.style={enhanced,breakable,colback=pale,colframe=teal!45!white,boxrule=.6pt,arc=1mm,left=3mm,right=3mm,top=2mm,bottom=2mm}}
\newcommand{\LESSONNUM}{30}
\begin{document}

{\Huge\sffamily\bfseries\color{deep} Kombinatorické pravidlo součtu a součinu}\par
{\sffamily\large\color{teal} Výukový podklad · 30. vyučovací hodina}\par
\textbf{Hodinová dotace:} 45 minut\quad \textbf{Ročník:} 4. SPŠ\quad \textbf{Rozpis:} TP 2.1\par
\begin{tcolorbox}[theobox]\textbf{Výukový cíl.} Používat pravidlo součtu pro disjunktní případy a pravidlo součinu pro postupné volby, včetně omezení na jednotlivých pozicích.\end{tcolorbox}
\sect{Teorie a vysvětlení}
\textbf{Pravidlo součtu.} Jestliže se výsledek uskuteční právě jedním ze vzájemně se vylučujících způsobů A nebo B a A má \(a\) možností, B má \(b\), celkem \(a+b\). Sčítat smíme jen \emph{nepřekrývající se} případy.\par
\textbf{Pravidlo součinu.} Jestliže první krok má \(a\) možností a po každé volbě prvního kroku následuje právě \(b\) možností druhého kroku, existuje \(ab\) dvojic. Při více krocích násobíme jednotlivé počty.\par
\textbf{Důležitá podmínka.} Jestliže počty pokračování závisejí na předchozí volbě, rozvětvíme řešení na případy a sečteme jejich výsledky. Stromový diagram pomáhá hlídat, zda jsme žádnou větev nevynechali.
\sect{Řešené příklady}
\textbf{Příklad 1.} Jízdenku lze koupit ve čtyřech různých tarifech na vlak nebo ve třech tarifech na autobus. Kolik je možností jedné jízdenky?\par
\textit{Řešení.} Volíme buď vlak, nebo autobus, nikoli oba zároveň. Množiny typů jízdenek jsou disjunktní. \(4+3=7\).\par\smallskip
\textbf{Příklad 2.} Na oběd je pět hlavních jídel a tři nápoje. Vybíráme jedno z každého.\par
\textit{Řešení.} Nejprve \(5\) možností jídla a pro každé z nich \(3\) možnosti nápoje. \(5\cdot3=15\).\par\smallskip
\textbf{Příklad 3.} Kolik existuje čtyřciferných čísel bez opakování číslic?\par
\textit{Řešení.} První cifra nemůže být \(0\): \(9\) možností. Druhá může být nula, ale nesmí být první cifrou: \(9\). Další cifry: \(8\) a \(7\). Počet \(9\cdot9\cdot8\cdot7=4536\).\par\smallskip
\textbf{Příklad 4.} Z číslic \(0,1,2,3,4\) vytvoř trojciferná čísla s různými číslicemi dělitelná pěti.\par
\textit{Řešení.} Takové číslo končí cifrou \(0\) (cifra \(5\) není k dispozici). První cifra má \(4\) možnosti, druhá \(3\); poslední je pevně \(0\). Počet \(4\cdot3=12\).\par\smallskip
\textbf{Příklad 5.} Čtyřznakový kód z \(A,B,C,D\), v němž nesmějí být stejné sousední znaky.\par
\textit{Řešení.} První místo \(4\) možnosti; na každém z dalších tří míst jsou \(3\) možnosti (vše kromě bezprostředně předchozího znaku). \(4\cdot3^3=108\).\par\smallskip
\clearpage\sect{Procvičování pro studenty}
\begin{enumerate}
\item Tři různé příchutě a čtyři polevy: jedna příchuť a jedna poleva.
\item Do finále se lze dostat jednou z pěti postupových cest A nebo jednou ze sedmi cest B; cesty se nepřekrývají.
\item Kolik čtyřmístných kódů z číslic \(0\) až \(5\) lze napsat, je-li opakování dovoleno?
\item Kolik trojciferných čísel lze napsat z číslic \(0\) až \(6\), bez opakování?
\item Kolik existuje pětimístných hesel ze tří znaků, v nichž žádné dva sousední nejsou stejné?
\end{enumerate}
\sect{Řešení a zdůvodnění k procvičování}
\begin{enumerate}
\item Součin \(3\cdot4=12\).
\item Součet \(5+7=12\).
\item Každé místo má \(6\) možností: \(6^4=1296\); kód může začínat nulou.
\item První cifra \(6\) možností, druhá \(6\), třetí \(5\): \(6\cdot6\cdot5=180\).
\item První znak \(3\), každý další \(2\): \(3\cdot2^4=48\).
\end{enumerate}
\sect{Metodické poznámky}
\textbf{Typické chyby:}
\begin{itemize}[leftmargin=5mm,itemsep=1pt]
\item Součet neznamená „udělám první volbu, potom druhou“; to je součin.
\item Nesmí se sčítat překrývající se případy bez korekce.
\item U čísla je první cifra nenulová, u PINu může být nula.
\end{itemize}
\textbf{Doporučený průběh 45 minut:} 0–7 navázání na výpis; 7–18 součet/součin a strom; 18–30 omezení číslic; 30–41 úlohy; 41–45 kontrola.\par
\textbf{Návaznost:} Na počítání dlouhých řetězců naváže faktoriál a zkracování součinů.\par
\vspace{2mm}\noindent{\footnotesize\color{slate} \textbf{Východiska:} rozpis výuky MAT, 4. ročník, hodiny 29--41; tematický plán SPŠ; srovnání tematického rozsahu a náročnosti s realisticky.cz, oddíl Kombinatorika (lekce 090101--090119). Výklad, příklady a řešení jsou nově autorsky sestavené, nikoli převzatá sbírka.}\par
\end{document}